
\medskip
\noindent\textit{Proof.}
\phantomsection
\label{sga3:VIIB:0.2.D-proof}
Let \((L_j)\) be the filtered decreasing family of open submodules of
\(M\). We endow \(M/K\) with the quotient topology; that is, a basis
of neighborhoods of \(0\) is formed by the open submodules
\((K+L_j)/K\). Since \(K\) is closed, it is equal to the intersection
of the \(K+L_j\), and hence the map
\[
  \tau:M/K\longrightarrow\varprojlim_j M/(K+L_j)
\]
is injective. It is also open, since the right-hand side is the
inverse limit of the discrete modules \(M/(K+L_j)\). Moreover, for
each \(j\), the map
\[
  t_j:M/L_j\longrightarrow M/(K+L_j)
\]
is surjective and has Artinian kernel. Thus, by the preceding lemma,
the map \(t\) in the following commutative diagram is surjective:
\begin{center}
  \phantomsection
  \label{sga3:VIIB:0.2.D-quotient-completion-square}
  \begin{tikzcd}[row sep=2.8em,column sep=4.6em]
M \arrow[r,"p","{\sim}"'] \arrow[d]
  & \varprojlim_j M/L_j \arrow[d,"t"] \\
M/K \arrow[r,"\tau"]
  & \varprojlim_j M/(K+L_j).
\end{tikzcd}

\end{center}
Since \(p\) is an isomorphism because \(M\) is complete, it follows
that \(\tau\) is surjective, and hence is an isomorphism. This proves
\textup{\SGARefLink{sga3:vii:paragraph:5-1:equation:i}{(i)}}.

\phantomsection
\label{sga3:VIIB:0.2.D-proof-ii-a}
Let us prove \textup{\SGARefLink{sga3:vii:paragraph:5-1:equation:ii}{(ii)}(a)}. By the foregoing, for every \(i\) there
is an isomorphism
\[
  M/M_i\xrightarrow{\sim}\varprojlim_j M/(M_i+L_j),
\]
so the two horizontal arrows in the following commutative diagram are
isomorphisms:
\begin{center}
  \phantomsection
  \label{sga3:VIIB:0.2.D-filtered-family-square}
  \begin{tikzcd}[row sep=2.8em,column sep=4.6em]
M \arrow[r,"p","{\sim}"'] \arrow[d,"g"]
  & \varprojlim_j M/L_j \arrow[d,"s"] \\
\varprojlim_i M/M_i \arrow[r,"{\sim}"']
  & \varprojlim_{i,j} M/(M_i+L_j).
\end{tikzcd}

\end{center}
Moreover, for each \(j\), the family of submodules
\((M_i+L_j)/L_j\) has a least element, since \(M/L_j\) is Artinian.
Consequently, the morphism
\[
  s_j:M/L_j\longrightarrow\varprojlim_i M/(L_j+M_i)
\]
is surjective; therefore, by the preceding lemma, \(s\) is
surjective. It follows that \(g\) is surjective. Finally, the kernel
of \(g\) is the inverse limit of the \(M_i\), that is, their
intersection. This proves \textup{\SGARefLink{sga3:VIIB:0.2.D:item:a}{(a)}}.

\phantomsection
\label{sga3:VIIB:0.2.D-proof-ii-b}
Let us deduce \textup{(b)} from this. Since \(N\) is a closed
submodule, and hence separated and complete, it is a pseudocompact
module for the topology induced by that of \(M\). Therefore, by
\textup{(a)}, the morphisms \(f\) and \(g\) in the following
commutative exact diagram are surjective:
\begin{center}
  \phantomsection
  \label{sga3:VIIB:0.2.D-snake-exact-diagram}
  \begin{tikzcd}[row sep=2.8em,column sep=2.15em]
0 \arrow[r] & N \arrow[r] \arrow[d,"f"]
  & M \arrow[r] \arrow[d,"g"]
  & M/N \arrow[r] \arrow[d,"h"]
  & 0 \\
0 \arrow[r] & \varprojlim_i N/(N\cap M_i) \arrow[r]
  & \varprojlim_i M/M_i \arrow[r]
  & \varprojlim_i M/(N+M_i)\, . & {}
\end{tikzcd}

\end{center}
The ``snake lemma'' then shows that the sequence
\[
  0\longrightarrow\ker f\longrightarrow\ker g
  \longrightarrow\ker h\longrightarrow 0
\]
is exact, and the equality
\[
  N+\bigcap_i M_i=\bigcap_i(N+M_i)
\]
follows.

\phantomsection
\label{sga3:VIIB:0.2.B-proof}
We can now prove \SGARefLink{sga3:vii:proposition:0-2-b}{Proposition~0.2.B}. Let \(f:M\to N\) be a morphism of
pseudocompact \(A\)-modules. Then \(K=\ker(f)\) is a closed submodule
of \(M\), and hence is separated and complete; thus \(K\), with the
topology induced by that of \(M\), is a pseudocompact module. By
\SGARefLink{sga3:VIIB:0.2.D}{0.2.D}\,\textup{\SGARefLink{sga3:vii:paragraph:5-1:equation:i}{(i)}}, the module \(M/K\), endowed with the quotient
topology, is pseudocompact.

Let us show that the continuous bijective morphism
\[
  M/K\longrightarrow\operatorname{Im}(f)
\]
is bicontinuous. After identifying \(M/K\) with
\(\operatorname{Im}(f)\), it is enough to show that the topology
\(T\) induced by that of \(N\) is finer than the quotient topology
\(Q\).\footnote{The French re-edition says here that \(Q\) is finer
than \(T\); the ensuing argument proves the reverse inclusion of open
sets, which is the one needed for bicontinuity.} Let \((L_j)\),
respectively \((N_i)\), be the filtered
decreasing families of open submodules of \(M\), respectively \(N\),
and put
\[
  N_i'=N_i\cap\operatorname{Im}(f).
\]
Let \(P=(K+L_j)/K\) be a submodule of \(M/K\) open for \(Q\). Since
\(M/(K+L_j)\) is Artinian, the family \(N_i'+P\) has a least element
\(N_{i_0}'+P\). The \(N_i'\) are open, and hence closed, for \(T\),
and therefore also closed for \(Q\). It follows from
\SGARefLink{sga3:VIIB:0.2.D}{0.2.D}\,\textup{\SGARefLink{sga3:vii:paragraph:5-1:equation:ii}{(ii)}\SGARefLink{sga3:VIIB:0.2.D:item:b}{(b)}} that
\[
  N_{i_0}'+P
  =\bigcap_i(N_i'+P)
  =P+\bigcap_i N_i'
  =P,
\]
whence \(N_{i_0}'\subset P\). Thus \(P\) is open for \(T\), and
\(M/K\to\operatorname{Im}(f)\) is consequently an isomorphism of
pseudocompact modules.

In particular, \(\operatorname{Im}(f)\) is complete for \(T\), and
hence closed in \(N\). By \SGARefLink{sga3:VIIB:0.2.D}{0.2.D}\,\textup{\SGARefLink{sga3:vii:paragraph:5-1:equation:i}{(i)}} once more,
\(\operatorname{Coker}(f)\), endowed with the quotient topology, is
pseudocompact. This proves that \(\operatorname{PC}(A)\) is an
abelian category.

Arbitrary inverse limits also exist in \(\operatorname{PC}(A)\): if
\((M_i)\) is an inverse system of pseudocompact modules, its inverse
limit has as
underlying module the inverse limit of the underlying
modules, endowed with the inverse-limit topology. Moreover, if one
has a family of exact sequences in \(\operatorname{PC}(A)\),
\[
  0\longrightarrow K_i\longrightarrow M_i
  \longrightarrow Q_i\longrightarrow 0,
\]
then the sequence
\[
  0\longrightarrow\prod_i K_i\longrightarrow\prod_i M_i
  \longrightarrow\prod_i Q_i\longrightarrow 0
\]
is exact. Assertion \textup{\SGARefLink{sga3:vii:proposition:0-2-b:item:iii}{(iii)}} of \SGARefLink{sga3:vii:proposition:0-2-b}{0.2.B} follows, since in every
abelian category in which arbitrary products exist, conditions
\textup{\SGARefLink{sga3:VIIB:0.2.D:item:a}{(a)}} and \textup{\SGARefLink{sga3:VIIB:0.2.D:item:b}{(b)}} of \SGARefLink{sga3:VIIB:0.2.D}{0.2.D} are equivalent, and are
equivalent to the exactness of filtered inverse limits
(\textup{cf.}\ [Mi65], III, \SGARefLink{sga3:iii:subsection:1-2}{1.2}--1.9, or [Po73], Chap.~2,
Th.~8.6).

Finally, every pseudocompact module \(M\) is a submodule of the
product
\[
  \prod_L M/L,
\]
where \(L\) ranges over the open submodules of \(M\). Thus the objects
of finite length form a system of cogenerators of
\(\operatorname{PC}(A)\). (Moreover, every object of length \(n\) is
isomorphic to a quotient \(A^n/L\), where \(L\) is an open submodule
of \(A^n\) of colength \(n\); these quotients therefore form a
\emph{set} of cogenerators.) This completes the proof of \SGARefLink{sga3:vii:proposition:0-2-b}{0.2.B}.

